% Generated draft: requires mathematical and source review.
\documentclass[11pt]{article}
\usepackage[margin=25mm]{geometry}
\usepackage{fontspec}
\setmainfont{Latin Modern Roman}
\usepackage{amsmath,amssymb,mathtools,unicode-math}
\setmathfont{Latin Modern Math}
\usepackage{enumitem,needspace,xurl,hyperref}
\setlength{\parindent}{0pt}
\setlength{\parskip}{5pt}
\setlength{\emergencystretch}{3em}
\newcommand{\N}{\mathbb{N}}
\newcommand{\Z}{\mathbb{Z}}
\newcommand{\Q}{\mathbb{Q}}
\newcommand{\R}{\mathbb{R}}
\newcommand{\C}{\mathbb{C}}
\newcommand{\F}{\mathbb{F}}
\newcommand{\A}{\mathbb{A}}
\newcommand{\PP}{\mathbb{P}}
\newcommand{\E}{\mathbb{E}}
\newcommand{\eps}{\varepsilon}
\newcommand{\dd}{\,\mathrm{d}}
\begin{document}
\section*{UnsolvedMath ID: 701627}
\textbf{Source problem code:} catalogue entry
\subsection{Definitions and mathematical statement}
Let \(G\) be an abstract group. For each integer \(n\geq 1\), a complex linear representation of dimension \(n\) is a group homomorphism \(\rho\colon G\to\operatorname{GL}_n(\C)\). It is irreducible if the only \(\rho(G)\)-invariant complex linear subspaces of \(\C^n\) are \(\{0\}\) and \(\C^n\). Two such representations \(\rho\) and \(\sigma\) are isomorphic if there exists \(A\in\operatorname{GL}_n(\C)\) such that \(\sigma(g)=A\rho(g)A^{-1}\) for every \(g\in G\).

Suppose that, for every integer \(n\geq 1\), there exists exactly one isomorphism class of irreducible complex linear representations of \(G\) of dimension \(n\). Must there exist a normal subgroup \(N\trianglelefteq G\) such that
\[
G/N\cong\operatorname{SL}_2(\Q),
\qquad
\operatorname{SL}_2(\Q)=\{M\in\operatorname{Mat}_{2\times2}(\Q):\det M=1\}?
\]
The supplied source states that \(G=\operatorname{SL}_2(\Q)\) satisfies the representation hypothesis, citing a theorem of Borel--Tits. This example and its cited justification are retained as supplied source claims and are not independently verified here. The record names J.-P. Serre as proposer and prefixes the statement with ``(T. Springer)''; the relationship between these attributions is not explained.
\subsection{Short English statement}
If an abstract group has exactly one irreducible complex representation, up to isomorphism, in every positive finite dimension, must it have a quotient isomorphic to \(\operatorname{SL}_2(\Q)\)? The supplied registry labels this problem open; that label has not been independently verified.
\subsection{Sources}
\begin{itemize}
\item Supplied local UnsolvedMath registry record, permanent ID \(701627\), exact source code \texttt{catalogue entry}, titled \emph{Kourovka Notebook Problem 19.79}. The record labels the problem open and was updated on 2026-10-16; this is supplied registry metadata, not an independently verified literature assessment.
\item \emph{The Kourovka Notebook (Unsolved Problems in Group Theory)}, original problem 19.79, 19th issue (2018); the supplied background also identifies the 21st issue edition (2026), arXiv:1401.0300v48. The supplied collection description names E. I. Khukhro and V. D. Mazurov as editors.
\item Attribution supplied by the registry: J.-P. Serre, proposed in 2018. The statement itself begins with the additional attribution ``(T. Springer)'', whose precise role is unspecified.
\item Borel--Tits, \emph{Ann. Math.}, volume \(97\) (1973), pages 499--571, cited in the supplied statement for the assertion that \(\operatorname{SL}_2(\Q)\) has the stated representation property. No theorem number is supplied, and this citation has not been independently checked.
\end{itemize}
\end{document}
